\section{导读：机器人数据荒到底荒在哪里}

本节先建立整期的阅读方式。EP109 不是泛泛地说“机器人需要数据”，而是把数据荒拆成四层：真实机器人数量不足，物理交互难采集，仿真和真实之间有 gap，以及数据是否真正提升模型还需要评价系统验证。谢晨的经历横跨 Cruise、NVIDIA、蔚来和光轮智能，因此他的叙述有一个很实用的特点：每次谈技术概念，都会回到“它能不能让模型变好、能不能让系统落地、能不能形成商业闭环”。

这期访谈也给张小珺机器人系列补上一个关键拼图。EP121 从 DeepMind 视角讲机器人基座模型、跨本体和世界模型；EP132 从星海图视角讲整机、供应链和 Data Recipe；EP109 则把底层的数据生产系统讲透：为什么合成数据不是“偷懒替代真实数据”，而是一套必须被真实反馈、评价系统、物理参数和客户需求持续校准的基础设施。

\begin{importantbox}{本期核心命题}
具身智能的数据荒不是“视频太少”或“机器人太少”这么简单，而是缺少可规模化、可验证、可服务强化学习和真实部署的物理世界数据。合成数据与仿真的价值，取决于它们能否构成 Real2Sim、Simulation、Sim2Real 和评价系统之间的闭环。
\end{importantbox}

讲者强调，一个看起来很炫的仿真世界，如果不能让感知模型、预测模型或机器人策略变好，只是一个漂亮玩具。真正的数据基础设施要从机器学习目标反推：哪些场景稀缺，哪些变量影响模型，哪些合成样本能带来真实 metric 提升，哪些 gap 必须被物理采集和真实部署反馈修正。

\begin{figure}[H]
\centering
\includegraphics[width=0.96\textwidth]{sim2real-loop.png}
\caption{Sim2Real 闭环：从仿真生成数据，到现实部署，再用真实反馈修正仿真。自制概念图，依据 00:02:48--00:16:17 对谈内容整理。}
\end{figure}

\begin{knowledgebox}{读图：合成数据不是替代现实，而是把现实放大}
这张图的起点不是“凭空生成”，而是真实世界中的稀缺问题：行人突然窜出、雾天、坡道、抓取失败、冰箱门铰链等。仿真把这些问题参数化、组合化、规模化，训练后再回到真实世界验证。读图时要看闭环箭头：如果真实反馈不能回到仿真，合成数据会越做越像内部自嗨；如果评价系统不能衡量模型是否提升，数据量再大也无法证明有效。
\end{knowledgebox}

\begin{knowledgebox}{视觉策略说明}
本视频是固定访谈画面，没有讲义、白板或产品演示。正文遵循播客工作流：封面只放在首页，正文不重复使用访谈画面；所有正文图像都是自制概念图，用来承载 Sim2Real、Real2Sim、产业链和 scaling law 等教学结构。
\end{knowledgebox}

\subsection{本章小结}

EP109 的主线可以压缩成一句话：具身智能要 scale，不能只等真实机器人自然产生数据，而要建立一套能被现实校准、能被模型评价、能服务真实部署的合成数据和仿真基础设施。

